\chapter{Structure of the Proof}\label{ch:structure}

The proof of the main theorem requires overcoming two fundamental challenges: the \textbf{non-compactness} of the ambient manifold $M$, and the \textbf{presence of the boundary} $\partial M$. This chapter provides an overview of how these challenges shape the technical framework and guide the proof strategy.

%%%%%%%%%%%%%%%%%%%%%%%%%%%%%%%%%%%%%%%%%%%%%%%%%%%%%%%%%%%%%%%%%%%%%%%%%%%%%%%
\section{Challenges from Non-Compactness}\label{sec:challenges-noncompact}
%%%%%%%%%%%%%%%%%%%%%%%%%%%%%%%%%%%%%%%%%%%%%%%%%%%%%%%%%%%%%%%%%%%%%%%%%%%%%%%

In the compact setting, the classical Almgren--Pitts min-max theory produces minimal hypersurfaces by considering sweepouts of the entire manifold. When $M$ is complete but non-compact, we must fundamentally rethink the setup.

\subsection*{Localizing the Sweepout}

Since $M$ is non-compact, the width of sweepouts of the entire manifold would typically be infinite or ill-defined. Following Chambers--Liokumovich~\cite{CL20}, we instead consider \emph{sweepouts of bounded open sets}: given a bounded open set $U \subset M$, a sweepout is a continuous family of relative cycles $\{\tau_t\}$ that starts ``outside $U$'' and ends having ``filled $U$.'' The width $W(U)$ is defined as the infimum of the maximal slice mass over all such sweepouts.

This localization introduces two difficulties: \textbf{stationarity} (how to perform pull-tight on non-compact $M$) and \textbf{escape} (how to ensure mass doesn't escape from $U$).

\subsection*{Stationarity and Escape}

These are solved by two independent mechanisms:

\begin{itemize}
    \item \textbf{Good set condition solves escape.} We require $\mathbf{M}(\partial_I U) \leq \frac{1}{10} W_\partial(U)$. This makes $\partial_I U$ a ``bottleneck'' too small to transfer the filling alone, forcing every good sweepout to have critical slices with substantial mass inside $U$.

    \item \textbf{Exhaustion $+$ pull-tight solves stationarity.} Since $M$ is non-compact, we pull-tight layer by layer on an exhaustion $U_0 \Subset U_1 \Subset \cdots \Subset M$, then use a diagonal argument to obtain a globally stationary varifold.
\end{itemize}

The key point is that these two mechanisms must work together: we need the stationary slice to come from a good sweepout of a good set. Pull-tight preserves good sweepout structure (since endpoints have small mass and are not modified), so the extracted stationary slice inherits the non-escape property from the good set condition.
